To begin, enable use of the \OHPC{} repository by adding it to the local list
of available package repositories. Note that this requires network access from
your {\em master} server to the \OHPC{} repository, or alternatively, that
the \OHPC{} repository be mirrored locally. For Ubuntu, the repository is a
signed APT repository and the example which follows installs the repository
key and source list from the VersatusHPC mirror.

% begin_ohpc_run
% ohpc_comment_header Install OpenHPC ohpc-release \ref{sec:enable_repo}
\begin{lstlisting}[language=bash,keywords={},basicstyle=\fontencoding{T1}\fontsize{7.6}{10}\ttfamily,
    literate={-}{-}1]
[sms](*\#*) export OHPC_REPO_ROOT=https://repos.versatushpc.com.br/openhpc/versatushpc-4
[sms](*\#*) export OHPC_UBUNTU_REPO=${OHPC_REPO_ROOT}/Ubuntu_24.04
[sms](*\#*) install -d -m 0755 /usr/share/keyrings /etc/apt/sources.list.d
[sms](*\#*) curl -fsSL ${OHPC_REPO_ROOT}/versatushpc.gpg \
             -o /usr/share/keyrings/versatushpc.gpg
[sms](*\#*) curl -fsSL ${OHPC_UBUNTU_REPO}/versatushpc-openhpc.list \
             -o /etc/apt/sources.list.d/versatushpc-openhpc.list
[sms](*\#*) apt update
[sms](*\#*) (*\install*) ohpc-release
\end{lstlisting}
% end_ohpc_run

\begin{center}
\begin{tcolorbox}[]
\small Many sites may find it useful or necessary to maintain a local copy of the
\OHPC{} repositories. For Ubuntu, mirror the signed APT repository rooted at
\href{https://repos.versatushpc.com.br/openhpc/versatushpc-4/Ubuntu_24.04/}
{\color{blue}{https://repos.versatushpc.com.br/openhpc/versatushpc-4/Ubuntu\_24.04/}}
and update \texttt{/etc/apt/sources.list.d/versatushpc-openhpc.list} to point at
the local mirror.
\end{tcolorbox}
\end{center}
